% META: {"id":"1NSI-TC-EX-049","chapitre":"1NSI-TYPES-CONSTRUITS","type_objet":"exercice","capacites":["1NSI-TYPES-CONSTRUITS-C5"],"parcours":1,"competences":["analyser"],"duree_min":8,"mode_creation":"ex_nihilo","sources_inspiration":[],"status":"approved","origin":"needs_review","status_history":["needs_review","audited","accepted","approved"],"release_acceptance":"RELEASE_OWNER_BATCH_ACCEPTANCE","acceptance_closure_digest":"sha256:9b3ccf9a81c5520fbb7e03b7b2d3e7bf2057a02834b6d908bf3be5f49f3b553f","acceptance_receipt_digest":"sha256:4fad86f0282e0fa4e0419cca1ad6379ae7af5a280c04a4a90880f90a1c044242"}
% BEGIN-TRACE
% t = (1, 2, 3)
% try:
%     t[0] = 10
% except TypeError:
%     print("erreur")
% print(t)
% EXPECTED
% erreur
% (1, 2, 3)
% END-TRACE
\begin{exercice}{1NSI-TC-EX-049}{1}{8}
On considère le programme suivant :
\begin{python}
t = (1, 2, 3)
t[0] = 10
print(t)
\end{python}
\begin{enumerate}
  \item Que se passe-t-il lorsqu'on exécute ce programme ? Quelle erreur obtient-on ?
  \item Expliquer pourquoi on dit qu'un tuple est \textbf{immuable} (non mutable).
  \item Donner un avantage d'utiliser un tuple plutôt qu'une liste quand on
        ne souhaite pas modifier les données.
\end{enumerate}
\end{exercice}
